\subsection{A smaller separating multiplier}
Put $B_*=P(P+R)$. The source bounds imply $U\ge64$, $M\ge512$, and
$A\ge33280$. In particular all the following inequalities hold for
\emph{every} assignment in the specified domain, before imposing any
square or divisibility condition:
\begin{equation}\label{eq:rootbounds}
 0<\sqrt{a_0}<PK\le hB_*<h\sqrt{a_1}.
\end{equation}
Indeed, $h(P+R)-K=(h-1)(R+1)\ge0$, and
$a_0=P^2K^2-K^2+1<P^2K^2$. Also
\[
 A^2-2P=M^2(U-1)^2>0,\qquad D>A^2>2P,\qquad E>D.
\]
It follows that $F>2P^2>B_*$, since $R<P$.
The quantity $L=CD+1-A$ is greater than one, so $G>F>B_*$.
Finally $I\ge G^2$ and $DF>1$, proving the last bound in
\eqref{eq:rootbounds}. Thus the existing coordinate $h$ already separates
the two radicals. No new unknown is needed for that purpose.

\subsection{The principal polynomial majorant}
Define $\lambda=2048$ and the following polynomial registers:
\begin{align*}
 L&=CD+1-A,& \mathcal N&=2G+L,\\
 \mathcal N_1&=2e^2D(D-1)L+3C,&
 \mathcal N_2&=2e^2(D-1)L,\\
 \mathcal N_3&=2e^2(L-C),& \mathcal N_4&=2e^2C,\\
 c_*&=256i(e^2+2)+2048,& B_g&=2P(P+R)+1.
\end{align*}
Put
\begin{align}
 T={}&e\bigl((1024D^2-512D-128)\mathcal N
       -64\mathcal N_1-40\mathcal N_2
       -28\mathcal N_3-21\mathcal N_4\bigr)+c_*,\label{eq:Tmajor}\\
 V={}&TH+B_g,\qquad b_1=\lambda^2a_1.\label{eq:Vmajor}
\end{align}
These definitions have no division. The registers $\mathcal N_k$ are
explicit products and sums, not quotient operations.

\begin{lemma}[Uniform polynomial majorant]\label{lem:majorant}
For every allowed assignment,
\[
 T>\lambda\sqrt{DF}\,G>0,
 \qquad \sqrt{b_1}<V-B_*,
 \qquad B_*<\sqrt{b_1}.
\]
\end{lemma}
\begin{proof}
It is enough for the first inequality that $A\ge3$ and $C,i\ge1$.
Let $X_*=e^2D^3(D-1)$. Direct substitution gives
\[
 DF=X_*+D,\qquad G=(X_*/D)L+CD+1.
\]
The fractions in this proof are real-number comparisons only.
Define the two integer polynomials
\begin{align*}
 t_6(D)={}&1024D^6-512D^5-128D^4-64D^3-40D^2-28D-21,\\
 d_5(D)={}&33792D^5+10560D^4+4928D^3+2464D^2+1176D+441.
\end{align*}
Their exact identity is
\begin{equation}\label{eq:tnidentity}
 t_6(D)^2-1024^2D^{11}(D-1)=d_5(D).
\end{equation}
For $D\ge1$, $t_6(D)>0$: the sum of the negative coefficients is 793,
less than 1024, and every corresponding power is at most $D^6$.
Also $d_5(D)>0$. Hence, for $D>1$,
\[
 \frac{t_6(D)}{D^4}>1024D^2\sqrt{1-1/D}.
\]
This finite identity replaces any need to assume convergence of a series.

The exact remainders $\mathcal N=D^k\mathcal N_k+r_k$ are
\begin{align*}
 r_1&=3-A, &r_2&=3CD+3-A,\\
 r_3&=2e^2(A-1)D^2+r_2, &
 r_4&=2e^2(1-A-C)D^3+r_3.
\end{align*}
They give the identity
\[
 T=e\frac{t_6(D)}{D^4}\mathcal N+\mathcal E+c_*,
 \qquad
 \mathcal E=e\left(\frac{64r_1}{D}+\frac{40r_2}{D^2}
                    +\frac{28r_3}{D^3}+\frac{21r_4}{D^4}\right).
\]
Since $D\ge A^2\ge9$ and $e=i(4A-5)C$, elementary bounds give
\[
 \frac{eA}{D}<\frac92 i,\qquad
 \frac{eC}{D}<\frac32 i,\qquad
 \frac{e(A+C-1)}{D}<6i.
\]
For example, the first follows from
$eA/D<4iA^2/((A^2-1)C)\le9i/2$; the second follows from
$eC/D<4iA/(A^2-1)\le3i/2$. Both $r_2$ and $r_3$ are positive.
Bounding only the possibly negative terms therefore yields
\begin{align*}
 -\mathcal E
 &<\frac{64eA}{D}+\frac{42e^3(A+C-1)}{D}+\frac{21eA}{D^4}\\
 &<252ie^2+\frac{765}{2}i<c_*-2048.
\end{align*}

For positive $X_*,D$, strict concavity of the square root gives
$\sqrt{X_*+D}<\sqrt{X_*}+D/(2\sqrt{X_*})$.
Multiplication by $2048G$ and the displayed identity for $G$ give
\[
 \lambda\sqrt{DF}\,G
 <1024eD^2\sqrt{1-1/D}\,(2G+L)+\lambda v_*,
 \quad
 v_*:=\frac{CD+1}{2eD\sqrt{1-1/D}}<1.
\]
The last inequality follows from $e\ge7C$ and
$\sqrt{D(D-1)}>D-1$. Combining these bounds proves
$T>\lambda\sqrt{DF}\,G$.
Since $I\le G^2H^2$, equations \eqref{eq:Vmajor} imply
$V-\sqrt{b_1}>B_g>B_*$. The remaining lower bound follows from
\eqref{eq:rootbounds}.
\end{proof}

The low-degree padding $B_g=2B_*+1$ has two roles. It supplies the strict
root gap for all inputs, and its factor two prevents the highest terms of
the final norm from cancelling. Omitting that factor preserves the gap
but invalidates the stated exact-degree calculation.

\subsection{The final equation and its exact existential meaning}
Set
\begin{align*}
 \Omega&=um,\\
 \Theta&=u(n+M_d)+v-\Omega(2M_d+hV).
\end{align*}
The principal nine-variable equation is
\begin{equation}\label{eq:ninefinal}
 \boxed{F_9=
 \bigl(\Theta^2-\Omega^2(a_0+h^2b_1)\bigr)^2
       -4\Omega^4h^2a_0b_1=0.}
\end{equation}
Together with the displayed source, register definitions, fixed coefficient
bytes, weights, and targets, this is a completely specified ordinary
integer polynomial in exactly nine unknowns.

\begin{proposition}[No lost sign or margin condition]\label{prop:ninenorm}
For fixed values of the first eight mixed-domain coordinates,
\eqref{eq:ninefinal} admits $n\in\N$ if and only if
$a_0,a_1$ are integer squares, $u\mid v$, and the computed margin $m$ is positive.
\end{proposition}
\begin{proof}
Here $u>0$ and $0<v/u<M_d-1$. The latter follows directly from
$v/u=Q/J+(H-C)/F$, with $Q>0$, $J,F>1$, and $H-C>0$.
If $m=0$, the left-hand polynomial is
$(u(n+M_d)+v)^4>0$. Thus a zero has $\Omega\ne0$.
Factoring the four-root norm over the real numbers gives
\[
 \frac\Theta\Omega=\varepsilon\sqrt{a_0}
                       +\delta h\sqrt{b_1},
 \qquad \varepsilon,\delta\in\{-1,1\}.
\]
This rational sum is nonzero by \eqref{eq:rootbounds}. If it is $r$, then
\[
 \sqrt{a_0}=\frac{r^2+a_0-h^2b_1}{2\varepsilon r}.
\]
Both roots are therefore rational and hence integral. This handles equal
or dependent square classes as well. Because $b_1=\lambda^2a_1$,
integrality of $\sqrt{b_1}$ implies integrality of $\sqrt{a_1}$: a rational
square root of an integer is an integer.

Write the signed integer roots as $r_0,r_1$ and define
\[
 Z=2M_d+hV+r_0+hr_1.
\]
The strict root bounds, with integer endpoints, give
$r_0\ge-hB_*+1$ and $r_1\ge-V+B_*+1$. Consequently
$Z\ge2M_d+h+1>0$. Rearranging the equation for $\Theta$ gives
\[
 n+M_d+v/u=mZ.
\]
Its left side is positive, so $m>0$. Its right side and $n+M_d$ are
integers, so $u\mid v$.

Conversely, suppose the three conditions hold. Choose either signs of
the integral roots, retain the \emph{actual} positive computed value of $m$,
and set $n=mZ-M_d-v/u$. The same bounds imply $n\ge0$, and substitution
makes \eqref{eq:ninefinal} vanish. No step replaces the computed margin
by a freely chosen witness or assumes $m=1$.
\end{proof}

Composing Proposition~\ref{prop:ninenorm} with the numerical ratio bridge
and the source lemma proves the nine-variable existence equivalence.

\subsection{A smaller-cost alternative}
The release also supplies the simpler choice $\lambda=8$,
\[
 S_8=e(8D^2-4D-1),\qquad V_8=S_8GH,\qquad b_1=64a_1.
\]
Use the same norm \eqref{eq:ninefinal} with $V=V_8$.
Its exact defect is
\[
 V_8^2-64a_1
 =\bigl(e^2(8D+1)-64D\bigr)G^2H^2+64DF(H^2-1).
\]
Since $e\ge7$, the first bracket is greater than $4De^2$.
Also $V_8<8D^2eGH$ and $2B_*<D^2$, so
\[
 \frac{V_8^2-64a_1}{V_8}
 >\frac{eGH}{2D}>\frac{(D-1)De^3}{2}>D^2>2B_*.
\]
As $V_8>B_*$, this implies $(V_8-B_*)^2>64a_1$ and proves the same
strict gap. The preceding norm proof applies unchanged.
This alternative uses 568 operations with uniform targets or 624 with
overlapping targets; its exact degrees are respectively 24,399,300 and
18,262,476. It is included because operation count and degree have
different minima among the supplied constructions.
